\documentclass[12pt,a4paper]{article}
\usepackage{amsmath,amssymb,amsfonts}
\usepackage[mathscr]{euscript}
\newcommand{\R}{\mathbb{R}}
\newcommand{\C}{\mathbb{C}}
\newcommand{\N}{\mathbb{N}}
\newcommand{\Z}{\mathbb{Z}}
\newcommand{\Q}{\mathbb{Q}}
\newcommand{\dist}{\mathcal{D}'}
\newcommand{\tors}{\operatorname{tors}}
\title{Generalised Greedy--Prime Correspondence for Automorphic \(L\)-functions}
\author{}
\date{}

\begin{document}
\maketitle

\begin{abstract}
We extend the Greedy--Prime Correspondence Conjecture from the Riemann zeta function to all principal automorphic \(L\)-functions.  For any such \(L\)-function \(L(s,\pi)\) with an unconditional explicit formula, we define a \emph{greedy harmonic set} relative to the arithmetic coefficients \(c(p^m)=\lambda_\pi(p^m)\) using a twisted harmonic series.  The conjecture identifies this set, after a suitable logarithmic rescaling, with the set of prime‑power logarithms, the weights being exactly the coefficients appearing in the prime‑sum side of the explicit formula.  If the conjecture holds, the associated geometric programme yields a proof of the generalised Riemann Hypothesis.
\end{abstract}

\section{Twisted Harmonic Series and the Greedy Algorithm}
Let \(L(s,\pi)\) be an (irreducible cuspidal) automorphic representation of \(\mathrm{GL}_n\) over \(\Q\), with standard \(L\)-function
\[
L(s,\pi)=\sum_{n=1}^{\infty}\frac{\lambda_\pi(n)}{n^{s}},\qquad \Re(s)>1,
\]
where \(\lambda_\pi(n)\) are the Hecke eigenvalues (with \(\lambda_\pi(1)=1\)).  The Euler product writes
\[
L(s,\pi)=\prod_{p}\prod_{j=1}^{n}\bigl(1-\alpha_{j,\pi}(p)p^{-s}\bigr)^{-1},
\]
and the coefficients satisfy a generalised von Mangoldt function \(\Lambda_\pi(p^m)=\sum_{j=1}^{n}\alpha_{j,\pi}(p)^m\log p\).  The explicit formula (Weil, Barner) becomes
\[
\sum_{\rho}\widehat{h}(\rho)+\text{arch}=\sum_{p,m}\frac{\Lambda_\pi(p^m)}{p^{m/2}}h(m\log p)+\text{similar sum for dual}.
\]

To adapt the greedy harmonic sieve, we consider the \emph{twisted harmonic series}
\[
H_\pi(n)=\sum_{k=1}^{n}\frac{\lambda_\pi(k)}{k}.
\]
When \(\lambda_\pi(k)\) is real (e.g.\ for self‑dual representations, where the coefficients are real), the partial sums \(H_\pi(n)\) are real numbers.  For complex coefficients one may take the real part, or construct a two‑dimensional greedy sweep; we restrict attention to real coefficients to keep the geometry one‑dimensional.

For \(\theta\in[0,\pi]\) we define a greedy set \(S_\pi(\theta)\subset\N_{>0}\) by the rule
\[
S_\pi(\theta)=\Bigl\{k\ge1:\text{ the partial sum } \sum_{j\in S_\pi(\theta),\,j\le k}\frac{\lambda_\pi(j)}{j}\le \frac{\theta}{\pi}\text{ is maximal}\Bigr\},
\]
constructed by scanning \(k=1,2,\dots\) and including \(k\) whenever the current residual \(\frac{\theta}{\pi}-\sum_{j\text{ already chosen}}\frac{\lambda_\pi(j)}{j}\) is at least \(\frac{\lambda_\pi(k)}{k}\) (assuming \(\lambda_\pi(k)>0\); if negative, one may include it only when it helps to reach the target – a more refined rule is needed).  To avoid technicalities, we assume \(\lambda_\pi(k)\ge0\) for all \(k\); this holds for example for the Riemann zeta function (\(\lambda_\pi(k)=1\) for all \(k\)) and for certain \(L\)-functions attached to elliptic curves with complex multiplication where the coefficients are non‑negative.  For general self‑dual representations we may take the absolute value or split into positive and negative parts separately; the essential conjecture remains the same.

The \emph{greedy harmonic set} \(\mathcal{G}_\pi\) is then the union of all \(S_\pi(\theta)\):
\[
\mathcal{G}_\pi=\bigcup_{\theta\in[0,\pi]} S_\pi(\theta),\qquad S_\pi(\theta)\subseteq S_\pi(\theta')\ \text{if}\ \theta\le\theta'.
\]

\section{Logarithmic Mapping}
Define positions via the twisted harmonic numbers
\[
u_k=\pi H_\pi(k-1)=\pi\sum_{j=1}^{k-1}\frac{\lambda_\pi(j)}{j},\qquad k\ge1,
\]
with \(u_1=0\).  The greedy harmonic set then yields a set of points \(\mathcal{U}_{\mathcal{G}_\pi}=\{u_k:k\in\mathcal{G}_\pi\}\).

\section{The Generalised Conjecture}
\begin{conjecture}[Generalised Greedy--Prime Correspondence]\label{conj:GGPC}
Let \(\pi\) be a self‑dual cuspidal automorphic representation whose Hecke eigenvalues \(\lambda_\pi(n)\) are non‑negative real numbers.  Then there exists an affine rescaling \(u\mapsto \alpha u+\beta\) (with \(\alpha,\beta\in\R\)) such that
\[
\alpha\mathcal{U}_{\mathcal{G}_\pi}+\beta=\{\,m\log p:p\text{ prime},\,m\ge1\text{ with }\lambda_\pi(p^m)>0\,\},
\]
and the harmonic weights \(\frac{\lambda_\pi(k)}{k}\) correspond to the von Mangoldt weights \(\frac{\Lambda_\pi(p^m)}{p^{m/2}}\) under this identification.
\end{conjecture}

For representations where the coefficients may vanish for some prime powers, the set on the right‑hand side is restricted to those \(p^m\) with non‑zero coefficient.  The conjecture states that the greedy sieving of the twisted harmonic series reproduces exactly the prime powers that are visible to the \(L\)-function, with the precise arithmetic weights.

\section{Consequences for the Generalised Riemann Hypothesis}
The geometric programme extends verbatim: the twin conical helices are built using the coefficients \(\lambda_\pi(n)\), the Möbius inversion (via the inverse of the \(L\)-function) yields the twisted prime distribution \(\tau_{\pi,\mathrm{prime}}(t)=\sum_{p,m}\frac{\Lambda_\pi(p^m)}{p^{m/2}}\delta(t-m\log p)\).  The Riemann helix for \(\pi\) is constructed from the imaginary parts of the zeros of \(L(s,\pi)\) with a phase function \(\phi_\pi\) satisfying \(\phi_\pi(\gamma_n)=2\pi n\).  The Geometric Explicit Formula asserts that the torsion of this helix equals \(\tau_{\pi,\mathrm{prime}}\) plus a smooth term.  If Conjecture~\ref{conj:GGPC} holds, one can define \(\phi_\pi\) via the greedy harmonic set as in the \(\zeta\) case, thereby establishing the GEF and, through the Dirac operator with point interactions, proving the Generalised Riemann Hypothesis (GRH) for \(L(s,\pi)\).

The conjecture is equivalent to GRH for the given representation; hence it is a faithful geometric reformulation.

\section{Examples and Remarks}
\begin{itemize}
\item For the Riemann zeta function, \(\lambda_\pi(n)=1\) and the conjecture reduces to the classical Greedy--Prime Correspondence.
\item For a Dirichlet \(L\)-function \(L(s,\chi)\) with a real primitive character \(\chi\) (e.g.\ Legendre symbol), \(\lambda_\chi(n)=\chi(n)\in\{0,\pm1\}\).  One may consider the positive and negative parts separately and form a pair of helical curves, the torsion of which captures the twisted prime distribution.
\item For modular form \(L\)-functions (e.g.\ associated to an elliptic curve), the coefficients \(\lambda_E(p)=a_p\) are the Frobenius traces, which are real (by the functional equation) but may be positive or negative.  The greedy sieve can be applied to the sequence of positive coefficients and negative coefficients independently, yielding two interwoven helices.
\end{itemize}

\section{Outlook}
The Generalised Greedy--Prime Correspondence Conjecture for all principal \(L\)-functions presents a uniform combinatorial path to the Grand Riemann Hypothesis.  Its proof—or disproof—would constitute a major breakthrough.  The geometric framework ensures that the conjecture is equivalent to the corresponding Riemann Hypothesis, thus any progress on the greedy sieve directly translates to progress on GRH.

\end{document}